% !TEX root = ../../../main.tex
% Block card: detection.deconvolve -- "Deconvolve (OASIS)" (Part I, the blocks chapter).
\cardfield{What it does}
Explains the trace as a list of sudden jumps, each followed by the same slow fade, and gives
every moment a score: how big a jump starts there, in units of the background noise. Almost
every score is zero.

\cardfield{When to use it}
When your events are troughs: a sudden drop, then a slow recovery. Its advantage over a
plain level cut is that a second drop starting before the first has recovered counts as an
event of its own instead of merging into one blob. Put a long \blockref{preprocessing.detrend}
in front and \blockref{detection.impulse_events} after it; the ready-made template
\texttt{deconvolved\_troughs} does both.

\begin{plainterms}[Jumps and fades]
Think of a bath with the plug out, slowly draining at the same rate every time. Each time
someone pours in a bucket, the level jumps and then drains away the same way. Looking only at
the water level, this block works out when buckets were poured and how big each was, even when
a second bucket went in before the first had drained. A fungal trough is the same picture
upside down: a sudden drop, then a slow climb back.
\end{plainterms}

\cardfield{What goes in, what comes out}
In goes a stretch of one channel. Out comes a score for every sample: zero where nothing
starts, and the size of the jump, counted in noise widths, where one does.

\cardfield{Settings that matter}
\emph{recovery\_half\_s} first: the time a trough takes to climb half-way back. Read it off a
few of the recording's own troughs. At 0 the block estimates it, as a fallback. A value much
shorter than the real recovery splits every trough into a jump and many small follow-ons.
Then \emph{min\_jump\_sigma}: the smallest jump allowed, in noise widths; raise it to stop
noise being explained as tiny jumps. \emph{baseline\_window\_s} lets the floor under the trace
follow slow drift; make it about 15 recovery half-times.

\cardfield{Watch out for}
Events that are not a sudden drop and the same fade (a sharkfin's slow climb and plateau) are
split into many small events. The summary flags this when the leftover misfit (the residual
SD) is above about 1.5: then the event count is not a trough count. The summary also warns
when jumps come faster than they could recover, which usually means the recovery half-time is
wrong. The scores do not depend on whether the recording is stored in volts or millivolts.
